\documentclass[12pt]{article}
\usepackage[landscape, left=.5in, right=.5in, top=.5in, bottom=.75in]{geometry}
\usepackage{amsmath,amssymb,amsfonts,amsthm}
\usepackage{arydshln}
\setlength{\dashlinegap}{1pt}
\setlength{\dashlinedash}{1pt}



\title{UF3 KOKKOS Tensor Core Implementation}
\author{Mitch Murphy}
\date{\today}

\begin{document}
\maketitle

%\tableofcontents\clearpage


\section*{Matrix and Tensor Operations for the UF3 Potential}

The provided C++ code evaluates an empirical potential (likely a machine-learning model like UF3) utilizing cubic B-splines for 2-body and 3-body interactions. 

\subsection*{1. Distance Polynomial Vectors}
First, we define the polynomial distance vectors used throughout the code to evaluate the splines. For any pair distance $r$, we define the $4 \times 1$ position vector $\mathbf{p}(r)$ and the $3 \times 1$ derivative vector $\mathbf{p}'(r)$:

\begin{equation}
\mathbf{p}(r) = \begin{bmatrix} 1 \\ r \\ r^2 \\ r^3 \end{bmatrix}, \quad 
\mathbf{p}'(r) = \begin{bmatrix} 1 \\ r \\ r^2 \end{bmatrix}
\end{equation}

\subsection*{2. Two-Body Interactions}
The 2-body scalar force and energy are computed using a set of spline coefficients. Let $C_{2b}$ be the $4 \times 4$ coefficient matrix (\texttt{constants\_2b}) and $C'_{2b}$ be the $3 \times 3$ derivative coefficient matrix (\texttt{constants\_2b\_deri}).

\textbf{2-Body Energy:}
\begin{equation}
E_{2b} = \mathbf{1}_{4}^T C_{2b} \mathbf{p}(r_{ij})
\end{equation}
\textit{(where $\mathbf{1}_{4}$ is a $4 \times 1$ vector of ones).}

\textbf{2-Body Scalar Force:}
\begin{equation}
F_{2b} = \mathbf{1}_{3}^T C'_{2b} \mathbf{p}'(r_{ij})
\end{equation}

\textbf{2-Body Vector Force:}
Using the displacement vector $\Delta\mathbf{r}_{ij} = \mathbf{r}_i - \mathbf{r}_j$, the force applied to atom $i$ is:
\begin{equation}
\mathbf{f}_{i}^{(2b)} = -\frac{F_{2b}}{r_{ij}} \Delta\mathbf{r}_{ij}
\end{equation}

\subsection*{3. Three-Body Interactions: Basis Vectors}
For a triplet of atoms $(i, j, k)$, the code computes spline basis vectors for each pair distance ($r_{ij}, r_{ik}, r_{jk}$). Let $C_{ij}$ be the $4 \times 4$ coefficient matrix (\texttt{cc\_3b\_ij}) and $C'_{ij}$ be the $3 \times 3$ derivative matrix (\texttt{cc\_3b\_deri\_ij}). 

The $4 \times 1$ basis vectors ($\mathbf{b}$) and $3 \times 1$ derivative basis vectors ($\mathbf{b}'$) are obtained via matrix-vector multiplication:

\begin{align}
\mathbf{b}_{ij} &= C_{ij} \mathbf{p}(r_{ij}), & \mathbf{b}'_{ij} &= C'_{ij} \mathbf{p}'(r_{ij}) \\
\mathbf{b}_{ik} &= C_{ik} \mathbf{p}(r_{ik}), & \mathbf{b}'_{ik} &= C'_{ik} \mathbf{p}'(r_{ik}) \\
\mathbf{b}_{jk} &= C_{jk} \mathbf{p}(r_{jk}), & \mathbf{b}'_{jk} &= C'_{jk} \mathbf{p}'(r_{jk})
\end{align}

\subsection*{4. Three-Body Interactions: Tensor Contractions}
The nested loops representing \texttt{triangle\_eval} variables are tensor contractions. The coefficients are stored in 3-dimensional tensors, which we can represent as sums of bilinear forms evaluated over matrix slices $W_l$.

\textbf{Three-Body Energy (\texttt{triangle\_eval0}):} \\
Let $W^{(E)}$ be the $4 \times 4 \times 4$ energy coefficient tensor (\texttt{n3b\_coeff\_array}). The energy is the contraction of the basis vectors along the three axes:
\begin{equation}
E_{3b} = \sum_{l=0}^{3} b_{ij, l} \left( \mathbf{b}_{ik}^T W^{(E)}_{l} \mathbf{b}_{jk} \right)
\end{equation}

\textbf{Three-Body Scalar Forces (\texttt{triangle\_eval1, 2, 3}):} \\
The code evaluates the derivative components by contracting the derivative basis vector of one pair with the standard basis vectors of the other two pairs.

For the $ij$ pair (\texttt{triangle\_eval1}), using the $4 \times 4$ matrix slices $W^{(ij)}_l$ from \texttt{coeff\_for\_der\_ij}:
\begin{equation}
T_1 = \sum_{l=0}^{2} b'_{ij, l} \left( \mathbf{b}_{ik}^T W^{(ij)}_{l} \mathbf{b}_{jk} \right)
\end{equation}

For the $ik$ pair (\texttt{triangle\_eval2}), using the $3 \times 4$ matrix slices $W^{(ik)}_l$ from \texttt{coeff\_for\_der\_ik}:
\begin{equation}
T_2 = \sum_{l=0}^{3} b_{ij, l} \left( (\mathbf{b}'_{ik})^T W^{(ik)}_{l} \mathbf{b}_{jk} \right)
\end{equation}

For the $jk$ pair (\texttt{triangle\_eval3}), using the $4 \times 3$ matrix slices $W^{(jk)}_l$ from \texttt{coeff\_for\_der\_jk}:
\begin{equation}
T_3 = \sum_{l=0}^{3} b_{ij, l} \left( \mathbf{b}_{ik}^T W^{(jk)}_{l} \mathbf{b}'_{jk} \right)
\end{equation}

\subsection*{5. Three-Body Vector Forces}
Finally, the scalar quantities $T_1$, $T_2$, and $T_3$ are projected back onto the cartesian unit vectors between the atoms. Let the unit displacement vectors be defined as $\hat{\mathbf{r}}_{ji} = \frac{\mathbf{r}_j - \mathbf{r}_i}{r_{ij}}$, $\hat{\mathbf{r}}_{ki} = \frac{\mathbf{r}_k - \mathbf{r}_i}{r_{ik}}$, and $\hat{\mathbf{r}}_{kj} = \frac{\mathbf{r}_k - \mathbf{r}_j}{r_{jk}}$.

The intermediate force vectors are:
\begin{align}
\mathbf{f}_{ij} &= T_1 \hat{\mathbf{r}}_{ji} \\
\mathbf{f}_{ik} &= T_2 \hat{\mathbf{r}}_{ki} \\
\mathbf{f}_{jk} &= T_3 \hat{\mathbf{r}}_{kj}
\end{align}

The total force distributed to atoms $i$, $j$, and $k$ from this triplet interaction is simply the linear combination of these intermediate vectors:
\begin{align}
\mathbf{F}_i &= \mathbf{f}_{ij} + \mathbf{f}_{ik} \\
\mathbf{F}_j &= -\mathbf{f}_{ij} + \mathbf{f}_{jk} \\
\mathbf{F}_k &= -\mathbf{f}_{ik} - \mathbf{f}_{jk}
\end{align}

\clearpage



\[
  \left[
    \begin{array}{@{} c c c c c @{}}
      1 & 2 & -1 &  3 &  4 \\
      \cdashline{1-1}
      0 & \multicolumn{1}{: c}{1} &  3 & -2 & -1 \\
      \cdashline{2-5}
      0 & 0 &  0 &  0 &  0 \\
      0 & 0 &  0 &  0 &  0
    \end{array}
  \right]
\]

\end{document}



\begin{equation}
\end{equation}



%\clearpage








\end{document}
